\documentclass[11pt]{article}
\usepackage{manual_docs_style}
\usepackage{booktabs}
\usepackage{tabularx}

\begin{document}

\docTitle{Claude Code Trajectory Logging and Context-Token Accounting}
\phantomsection\label{docstart:claude_code_context_accounting}

This note documents a limitation of historical TraceBench runs executed through
Claude Code: the archived Claude Code trajectory is a conversation and response
log, not a byte-for-byte capture of every model request.  The log is sufficient
to recover the visible dialogue, tool interactions, and provider-reported usage
totals, but it is not sufficient to independently reconstruct the complete model
context that produced those totals.

The limitation affects the interpretation of cumulative-input diagnostics in the
trajectory debug PDFs.  In particular, equality between an estimated count and a
provider count at the first measured agent step can be a consequence of
calibration rather than evidence that two independent measurements agree.


\section*{Recorded history versus model request}

The Claude Code history under
\code{agent_logs/.claude/projects/*.jsonl} records events emitted by the client.
Depending on the event, these records preserve the initial user prompt, subsequent
user messages, assistant content blocks, tool calls, tool results, selected
attachments, identifiers, timestamps, the Claude Code version, and provider usage
fields.

The model request is larger than this recorded history.  A request can also
contain a harness-generated system prefix, policy and environment instructions,
built-in tool definitions and their JSON schemas, dynamically injected context,
cache-control boundaries, and provider-specific request framing.  These
components affect the model input even when they are not emitted as ordinary
conversation events.

\begin{center}
\small
\begin{tabularx}{\textwidth}{@{}p{0.28\textwidth}p{0.20\textwidth}X@{}}
\toprule
\textbf{Component} & \textbf{Historical log} & \textbf{Consequence} \\
\midrule
Initial task prompt & Generally present & Its text can be retokenized. \\
Assistant messages and tool interactions & Generally present & The retained
conversation can be reconstructed subject to exporter normalization. \\
Selected harness attachments & Sometimes present & Only the emitted attachment
content is directly recoverable. \\
Provider usage counters & Present on measured responses & The provider total is
known, but the counters do not reveal missing text. \\
Complete system or harness prefix & Not captured as a complete request field & Its
historical text and ordering cannot be established from the trajectory alone. \\
Complete tool-definition array & Not captured & Tool calls reveal which tools were
used, but not necessarily the exact schemas supplied to the model. \\
Exact cache-block contents and boundaries & Not captured as a request snapshot & A
cache-read count is a size measurement, not a copy of the cached payload. \\
Exact provider request serialization & Not captured & Wrapper and formatting
effects cannot be replayed exactly from the trajectory alone. \\
\bottomrule
\end{tabularx}
\end{center}

Possessing the harness source or knowing the Claude Code version is useful but is
not equivalent to possessing the historical request.  Some prompt components are
generated dynamically from the working directory, configuration, available
tools, installed skills, environment state, and client version.  Recreating a
similar environment later does not prove that its serialized request is identical
to the request used by the archived run.


\section*{Claude provider usage fields}

For the Claude Code logs used here, the inclusive provider input count for one
model invocation is
\[
  P = T_{\mathrm{input}}
      + T_{\mathrm{cache\_creation}}
      + T_{\mathrm{cache\_read}}.
\]
The three terms correspond to the raw fields
\code{input_tokens}, \code{cache_creation_input_tokens}, and
\code{cache_read_input_tokens}.  They must be retained separately in source data
and summed exactly once when a single inclusive prompt count is required.

One observed first invocation reports:
\begin{DocCode}
input_tokens                 =     3
cache_creation_input_tokens  =  4898
cache_read_input_tokens      = 11113
                                 -----
inclusive prompt tokens      = 16014
\end{DocCode}
This establishes the provider-measured size of that input.  It does not supply the
text represented by all 16,014 tokens.  In particular, the 11,113-token cache-read
field states how many cached tokens were used, not what their content was.

Claude Code may emit one streamed assistant response as several JSONL records.
Those records can repeat the same provider message identifier and usage object.
An exporter must group fragments belonging to the same provider message, select
the final usage state for that message, and attach it once to the corresponding
canonical agent invocation.  Summing usage over raw fragments would count the same
invocation multiple times.


\section*{Why the missing prefix cannot be derived from its count}

A token count is not an invertible representation of text.  Many distinct system
prompts, tool schemas, and request layouts have the same token count.  Consequently,
the provider total and the reconstructed visible conversation determine only the
size of an omitted remainder under a chosen accounting rule; they do not determine
the remainder's text.

This matters especially when two tokenizers are involved.  Let
\(P_{\mathrm{Claude}}\) be the provider input count, and let
\(V_{\mathrm{o200k}}\) be the count obtained by applying
\code{o200k_base} to the visible retained content.  Defining
\[
  B = P_{\mathrm{Claude}} - V_{\mathrm{o200k}}
\]
does not reconstruct the omitted prefix under \code{o200k_base}.  The subtraction
mixes counts produced by different tokenization and accounting systems.  The result
\(B\) is a provider-calibrated offset, not a direct tokenization of the hidden
prefix.

If the subsequent displayed estimate is calculated as
\[
  \widehat{C} = B + V_{\mathrm{o200k}},
\]
then at the calibration step
\[
  \widehat{C}
  = (P_{\mathrm{Claude}} - V_{\mathrm{o200k}})
    + V_{\mathrm{o200k}}
  = P_{\mathrm{Claude}}.
\]
Thus a display such as
\begin{DocCode}
Cumulative input tokens: 16014 / 16014
\end{DocCode}
is exactly equal by construction when the first value uses this calibration.
It is not an independent agreement between \code{o200k_base} and Claude's token
accounting.


\section*{Consequences for historical TraceBench runs}

For an existing run without a complete request snapshot, TraceBench can report the
following facts honestly:
\begin{itemize}
  \item the provider input count for invocations whose usage was recorded;
  \item the provider output count for those invocations;
  \item an \code{o200k_base} count of the content that is actually present in the
  canonical trajectory;
  \item a provider-calibrated cumulative index, provided that it is explicitly
  labelled as calibrated and is not presented as an independent full-context
  estimate.
\end{itemize}

TraceBench cannot derive a pure \code{o200k_base} count of the complete historical
context when the hidden request components are absent.  Nor can it infer their text
from the provider usage counters.  Reconstructing visible messages and subtracting
their estimated count from the provider total does not remove this limitation.

The provider prompt count remains authoritative as a provider measurement.  The
limitation concerns only an independent reconstruction of the complete input and
the interpretation of comparisons against that provider measurement.


\section*{Recommended debug-PDF semantics}

The clearest representation for historical runs is to keep independently known
quantities separate:
\begin{itemize}
  \item \textbf{Provider cumulative input tokens}: display the inclusive provider
  count when available.
  \item \textbf{Reconstructed visible-context tokens}: display the
  \code{o200k_base} count of only the context components preserved in the
  trajectory.
  \item \textbf{Complete-context estimate}: display \,\code{-}\, unless the complete
  request was captured or the value is explicitly named a
  \emph{provider-calibrated estimate}.
\end{itemize}

If a two-value field is retained with the intended meaning
``independent complete-context estimate / provider measurement'', the strict
historical rendering is therefore:
\begin{DocCode}
Cumulative input tokens: - / 16014
\end{DocCode}
Alternatively, a calibrated value may be shown, but both the label and the manual
must make the dependency explicit.  The first calibrated pair must not be used as
validation evidence, because equality at the anchor is guaranteed algebraically.

Provider completion tokens and reconstructed step-added tokens also require
distinct labels.  A reconstructed step can include tool observations added by the
harness, whereas provider completion usage counts model output.  These quantities
need not cover the same content even before considering tokenizer differences.


\section*{Capture requirements for future runs}

An independently reproducible complete-context estimate requires a request-side
capture made immediately before each provider call.  The capture should preserve:
\begin{enumerate}
  \item the ordered system and harness instruction blocks;
  \item the complete ordered message list supplied to the model;
  \item every tool name, description, and input schema;
  \item attachments and dynamically injected environment context;
  \item cache-control markers, cache-block boundaries, and hashes of the exact
  cached content;
  \item model, client, harness, tokenizer, and schema versions;
  \item request options that can affect provider accounting;
  \item the exact UTF-8 text or a canonical serialization from which it can be
  reproduced; and
  \item the provider's raw usage fields, stored without discarding the cached and
  uncached split.
\end{enumerate}

The capture format should distinguish content from accounting metadata and should
record hashes for integrity checks.  Because system prompts, environment context,
and tool arguments can contain credentials or other sensitive data, capture and
publication must use an explicit redaction policy.  A redacted public artifact may
retain hashes and counts, while exact internal reproduction requires access to the
unredacted secured snapshot.

With such a snapshot, TraceBench can tokenize the same complete reconstructed
request with \code{o200k_base} for a model-independent diagnostic and compare it
with the provider measurement without using that measurement to manufacture the
estimate.  Without it, the visible history and the provider total must remain two
different kinds of evidence.

\end{document}
